% =============================================================================
%  Notas de clase — Sesión 17: MergeSort
%  Programación Avanzada — Universidad Panamericana
%  Compilar: pdflatex sesion17_mergesort.tex
% =============================================================================
\documentclass[12pt, oneside]{article}
\usepackage{progavanzada}

\begin{document}

\portada{Sesión 17}{MergeSort}{2026}

\tableofcontents
\newpage

% =============================================================================
\section{1945: ordenar 1\,000 tarjetas sin memoria suficiente}
% =============================================================================

En 1945, John von Neumann trabajaba en el diseño del EDVAC, uno de los primeros
computadores electrónicos. Se enfrentó a un problema concreto: necesitaba ordenar
listas de números que no cabían completas en la memoria del computador. La solución
que inventó: dividir la lista en partes manejables, ordenar cada parte por separado
en cinta magnética y luego \textit{mezclar} las cintas de dos en dos.

Ese algoritmo —\textbf{MergeSort}— tiene hoy 80 años y sigue siendo
la base del algoritmo de ordenamiento de Python, Java, y del ordenamiento estable
de la biblioteca estándar de C++.

\begin{tecnico}[Historia: von Neumann y el primer MergeSort]
  John von Neumann describió MergeSort en una nota interna de 1945 mientras
  trabajaba en el Institute for Advanced Study de Princeton. El contexto era
  puramente práctico: las cintas magnéticas podían leerse en secuencia pero no
  accederse al azar, así que el único algoritmo viable era uno que mezclara
  secuencias ya ordenadas.
  \par\smallskip
  La prueba de que Merge\-Sort es \textit{óptimo} vino décadas después,
  con la teoría de la información. Hoy sabemos que ningún algoritmo de
  ordenamiento basado en comparaciones puede ser asintóticamente más rápido.
  \par\smallskip
  \textit{Referencia:} Knuth, D.\,E. (1998). \textit{The Art of Computer
  Programming, Vol.\,3: Sorting and Searching}, 2.ª\,ed., \S5.2.4.
  Addison-Wesley.
\end{tecnico}

% =============================================================================
\section{La fila del salón}
% =============================================================================

Ocho estudiantes se paran en fila, cada uno con un número. Queremos ordenarlos.

\begin{center}
\begin{tabular}{|c|c|c|c|c|c|c|c|}
\hline
20 & 31 & 83 & 92 & 14 & 42 & 51 & 63 \\
\hline
\end{tabular}
\end{center}

La estrategia de MergeSort: dividir la fila por la mitad, ordenar cada mitad
\textit{por el mismo método} y luego unirlas. Aplicamos esto recursivamente hasta
llegar a filas de una sola persona — que ya están ordenadas por definición.

\begin{center}
\begin{tabular}{cccccccc}
  \multicolumn{8}{c}{\texttt{[20, 31, 83, 92, 14, 42, 51, 63]}} \\[4pt]
  \multicolumn{4}{c}{\texttt{[20, 31, 83, 92]}} &
  \multicolumn{4}{c}{\texttt{[14, 42, 51, 63]}} \\[4pt]
  \multicolumn{2}{c}{\texttt{[20, 31]}} &
  \multicolumn{2}{c}{\texttt{[83, 92]}} &
  \multicolumn{2}{c}{\texttt{[14, 42]}} &
  \multicolumn{2}{c}{\texttt{[51, 63]}} \\[4pt]
  \texttt{[20]} & \texttt{[31]} & \texttt{[83]} & \texttt{[92]} &
  \texttt{[14]} & \texttt{[42]} & \texttt{[51]} & \texttt{[63]} \\
\end{tabular}
\end{center}

Cada \texttt{[x]} es el caso base: una persona sola ya está ordenada. A partir
de ahí, subimos mezclando pares de filas ordenadas:

\begin{center}
\begin{tabular}{cccccccc}
  \texttt{[20]} & \texttt{[31]} & \texttt{[83]} & \texttt{[92]} &
  \texttt{[14]} & \texttt{[42]} & \texttt{[51]} & \texttt{[63]} \\[4pt]
  \multicolumn{2}{c}{\texttt{[20, 31]}} &
  \multicolumn{2}{c}{\texttt{[83, 92]}} &
  \multicolumn{2}{c}{\texttt{[14, 42]}} &
  \multicolumn{2}{c}{\texttt{[51, 63]}} \\[4pt]
  \multicolumn{4}{c}{\texttt{[20, 31, 83, 92]}} &
  \multicolumn{4}{c}{\texttt{[14, 42, 51, 63]}} \\[4pt]
  \multicolumn{8}{c}{\texttt{[14, 20, 31, 42, 51, 63, 83, 92]}} \\
\end{tabular}
\end{center}

% =============================================================================
\section{El algoritmo}
% =============================================================================

MergeSort tiene dos funciones que se llaman mutuamente:

\begin{center}
\begin{tabular}{lp{9cm}}
  \toprule
  \textbf{Función} & \textbf{Responsabilidad} \\
  \midrule
  \cpp{mergesort(arr, left, right)} &
    Divide el subarreglo en dos mitades, ordena cada una recursivamente,
    y llama a \cpp{merge} para unirlas. \\[4pt]
  \cpp{merge(arr, left, mid, right)} &
    Recibe dos mitades ya ordenadas y las fusiona en una sola
    secuencia ordenada, usando un buffer temporal. \\
  \bottomrule
\end{tabular}
\end{center}

\subsection{La función \texttt{mergesort}}

\begin{codigo}[caption={mergesort: divide y recurre}]
void mergesort(int* arr, int left, int right) {
    if (left >= right) return;            // caso base: 0 o 1 elemento

    int mid = (left + right) / 2;

    mergesort(arr, left,    mid);         // ordena mitad izquierda
    mergesort(arr, mid + 1, right);       // ordena mitad derecha
    merge    (arr, left, mid, right);     // fusiona las dos mitades
}
\end{codigo}

El punto medio se calcula como \cpp{(left + right) / 2}. La mitad izquierda
cubre los índices \cpp{[left, mid]} y la derecha \cpp{[mid+1, right]}.

\begin{recuerda}
  El \textbf{caso base} es \cpp{left >= right}, no solo \cpp{left == right}.
  Si \cpp{left > right} el subarreglo está vacío y tampoco hay nada que hacer.
  El \cpp{>=} lo cubre todo de forma segura.
\end{recuerda}

\subsection{La función \texttt{merge}}

Fusionar dos mitades ordenadas requiere un \textbf{buffer temporal}: no podemos
reubicar elementos en el arreglo original sin sobrescribir valores que todavía
necesitamos leer.

\begin{codigo}[caption={merge: fusión de dos mitades ordenadas}]
void merge(int* arr, int left, int mid, int right) {
    int n = right - left + 1;
    vector<int> temp(n);          // buffer temporal de tamano exacto

    int i = left, j = mid + 1, k = 0;

    // Compara el frente de cada mitad y toma el menor
    while (i <= mid && j <= right) {
        if (arr[i] <= arr[j])
            temp[k++] = arr[i++];
        else
            temp[k++] = arr[j++];
    }

    // Copia el resto de la mitad que no se agoto
    while (i <= mid)   temp[k++] = arr[i++];
    while (j <= right) temp[k++] = arr[j++];

    // Devuelve temp al arreglo original
    for (int k = 0; k < n; k++)
        arr[left + k] = temp[k];
}
\end{codigo}

\begin{ejemplo}[Traza de merge: \texttt{[20,31,83,92]} y \texttt{[14,42,51,63]}]
  \begin{center}
  \begin{tabular}{clll}
    \toprule
    \textbf{Paso} & \textbf{Comparación} & \textbf{Acción} & \textbf{temp} \\
    \midrule
    1 & $20 \leq 14$? No  & toma 14 & \texttt{[14]} \\
    2 & $20 \leq 42$? Sí  & toma 20 & \texttt{[14,20]} \\
    3 & $31 \leq 42$? Sí  & toma 31 & \texttt{[14,20,31]} \\
    4 & $83 \leq 42$? No  & toma 42 & \texttt{[14,20,31,42]} \\
    5 & $83 \leq 51$? No  & toma 51 & \texttt{[14,20,31,42,51]} \\
    6 & $83 \leq 63$? No  & toma 63 & \texttt{[14,20,31,42,51,63]} \\
    7 & Der. agotada        & copia 83,92 & \texttt{[14,20,31,42,51,63,83,92]} \\
    \bottomrule
  \end{tabular}
  \end{center}
\end{ejemplo}

\subsection{Paso por referencia: el contrato de la API}

Las firmas de ambas funciones son un ejemplo de diseño de API mediante paso por
referencia:

\begin{codigo}
void mergesort(int* arr, int left, int right)
void merge    (int* arr, int left, int mid, int right)
\end{codigo}

\begin{center}
\begin{tabular}{lll}
  \toprule
  \textbf{Parámetro} & \textbf{Tipo} & \textbf{Significado} \\
  \midrule
  \cpp{arr}   & \cpp{int*} & Referencia al arreglo — se modifica directamente \\
  \cpp{left}  & \cpp{int}  & Inicio del subarreglo actual (por valor: solo se lee) \\
  \cpp{right} & \cpp{int}  & Fin del subarreglo actual (por valor: solo se lee) \\
  \cpp{mid}   & \cpp{int}  & Punto de división (por valor: solo se lee) \\
  \bottomrule
\end{tabular}
\end{center}

En ningún momento se copia el arreglo. Todas las llamadas recursivas trabajan sobre
el \textbf{mismo bloque de memoria}. El tipo del parámetro comunica la intención:
\cpp{int*} significa ``voy a modificar esto''; \cpp{int} significa ``solo lo leo''.

% =============================================================================
\section{¿Cuánto tarda MergeSort?}
% =============================================================================

\subsection{Análisis informal}

El árbol de llamadas tiene $\log_2 n$ niveles. En cada nivel se procesan los
$n$ elementos exactamente una vez (repartidos entre las distintas llamadas a
\cpp{merge} de ese nivel). Total:

\[
  T(n) = n \cdot \log_2 n \quad \Rightarrow \quad \mathcal{O}(n \log n)
\]

A diferencia de Bubble Sort (\(\mathcal{O}(n^2)\)), MergeSort con $n=10^6$
elementos requiere \emph{veinte millones} de operaciones en lugar de
\emph{un billón}.

\begin{center}
\begin{tabular}{lrrrr}
  \toprule
  \textbf{Algoritmo} & $n=100$ & $n=10^3$ & $n=10^6$ & $n=10^9$ \\
  \midrule
  Bubble Sort \(\mathcal{O}(n^2)\)    &
    10\,000 & 1\,000\,000 & $10^{12}$ & $10^{18}$ \\
  MergeSort \(\mathcal{O}(n \log n)\) &
    700 & 10\,000 & $2\times 10^7$ & $3\times 10^{10}$ \\
  \bottomrule
\end{tabular}
\end{center}

\subsection{Complejidad en todos los casos}

MergeSort siempre divide por la mitad, sin importar los datos de entrada.
Por eso tiene la misma complejidad en el peor, promedio y mejor caso:
\(\mathcal{O}(n \log n)\) en los tres.

% =============================================================================
\section{¿Puede existir un algoritmo más rápido?}
% =============================================================================

Esta es una de las preguntas más hermosas de la informática teórica, y tiene
una respuesta definitiva: \textbf{no}. Ningún algoritmo de ordenamiento basado
en comparaciones puede ser asintóticamente más rápido que $\mathcal{O}(n \log n)$.

\subsection{El argumento de la información}

¿Cuántas formas distintas hay de ordenar $n$ elementos? La respuesta es $n!$
(\textit{n factorial}): para $n = 8$, son $8! = 40\,320$ posibles órdenes.

Cuando el algoritmo hace una comparación ($a \leq b$?), el resultado puede ser
verdadero o falso: separa el espacio de posibilidades en \textbf{dos mitades}.
Para distinguir entre $n!$ posibilidades usando preguntas de sí/no se necesitan
al menos:

\[
  \log_2(n!) \text{ comparaciones}
\]

Por la \textbf{aproximación de Stirling} ($\log_2(n!) \approx n \log_2 n - n
\log_2 e$):

\[
  \log_2(n!) = \Theta(n \log n)
\]

Cualquier algoritmo de ordenamiento por comparaciones necesita, en el peor caso,
$\Omega(n \log n)$ comparaciones. MergeSort usa exactamente $\Theta(n \log n)$:
\textbf{es óptimo}.

\begin{consejo}
  Hay algoritmos de ordenamiento más rápidos que \(\mathcal{O}(n \log n)\), como
  \textbf{Counting Sort} o \textbf{Radix Sort}, pero no funcionan por comparaciones.
  Explotan propiedades específicas de los datos (por ejemplo, que los números están
  acotados en un rango). Si los datos son enteros entre 0 y 1\,000, Counting Sort
  ordena en $\mathcal{O}(n)$ — pero falla si los datos son cadenas arbitrarias.
  MergeSort funciona con cualquier dato que pueda compararse.
\end{consejo}

\begin{tecnico}[El límite inferior fue probado por Knuth]
  David Knuth formalizó el límite inferior $\Omega(n \log n)$ para ordenamiento
  por comparaciones en el Volumen 3 de \textit{The Art of Computer Programming}
  (1973), usando árboles de decisión. La demostración completa usa la altura mínima
  de un árbol binario con al menos $n!$ hojas.
  \par\smallskip
  \textit{Referencia:} Knuth, D.\,E. (1998). \textit{TAOCP}, Vol.\,3, \S5.3.1.
\end{tecnico}

% =============================================================================
\section{Comparación con otros algoritmos}
% =============================================================================

\begin{center}
\begin{tabular}{lccccc}
  \toprule
  \textbf{Algoritmo} & \textbf{Mejor caso} & \textbf{Caso medio} &
  \textbf{Peor caso} & \textbf{Memoria extra} & \textbf{Estable} \\
  \midrule
  Bubble Sort    & $\mathcal{O}(n)$       & $\mathcal{O}(n^2)$     & $\mathcal{O}(n^2)$     & $\mathcal{O}(1)$      & Sí \\
  Selection Sort & $\mathcal{O}(n^2)$     & $\mathcal{O}(n^2)$     & $\mathcal{O}(n^2)$     & $\mathcal{O}(1)$      & No \\
  Insertion Sort & $\mathcal{O}(n)$       & $\mathcal{O}(n^2)$     & $\mathcal{O}(n^2)$     & $\mathcal{O}(1)$      & Sí \\
  \textbf{MergeSort}& $\mathcal{O}(n \log n)$ & $\mathcal{O}(n \log n)$ & $\mathcal{O}(n \log n)$ & $\mathcal{O}(n)$ & \textbf{Sí} \\
  QuickSort      & $\mathcal{O}(n \log n)$ & $\mathcal{O}(n \log n)$ & $\mathcal{O}(n^2)$    & $\mathcal{O}(\log n)$ & No \\
  HeapSort       & $\mathcal{O}(n \log n)$ & $\mathcal{O}(n \log n)$ & $\mathcal{O}(n \log n)$ & $\mathcal{O}(1)$    & No \\
  \bottomrule
\end{tabular}
\end{center}

\subsection{¿MergeSort es el mejor?}

\textbf{Depende del contexto.}

\begin{itemize}
  \item \textbf{Peor caso garantizado}: MergeSort y HeapSort son los únicos que
    garantizan $\mathcal{O}(n \log n)$ siempre. QuickSort puede degradarse a
    $\mathcal{O}(n^2)$ con entradas adversas (como un arreglo ya ordenado con
    pivote en el primer elemento).

  \item \textbf{Velocidad práctica en RAM}: QuickSort suele ser más rápido en
    la práctica porque accede a la memoria en patrones \textit{cache-friendly}
    (localidad de referencia). MergeSort salta entre la mitad izquierda y la
    derecha, causando más \textit{cache misses}.

  \item \textbf{Estabilidad}: MergeSort preserva el orden relativo de elementos
    iguales (si el alumno A y el alumno B tienen la misma calificación y A estaba
    antes que B, A sigue antes después de ordenar). Esto importa en bases de datos
    y aplicaciones donde el orden original tiene significado.

  \item \textbf{Ordenamiento externo}: cuando los datos no caben en RAM (archivos
    de terabytes), MergeSort es imbatible. Se cargan trozos, se ordenan en memoria,
    y se mezclan las secuencias resultantes en disco — exactamente lo que von Neumann
    resolvió con cintas magnéticas.
\end{itemize}

\begin{tecnico}[TimSort: el sucesor práctico de MergeSort]
  \textbf{TimSort} es el algoritmo que usa Python (\texttt{list.sort()}),
  Java (\texttt{Arrays.sort} para objetos) y la biblioteca estándar de C++
  (\texttt{std::stable\_sort}). Fue diseñado por Tim Peters en 2002.
  \par\smallskip
  La idea: en la práctica, los arreglos reales casi nunca son completamente
  aleatorios — tienen ``runs'' (subsecuencias ya ordenadas). TimSort detecta
  esos runs, los extiende con Insertion Sort hasta un tamaño mínimo, y los
  fusiona con el algoritmo de MergeSort. Obtiene $\mathcal{O}(n)$ en el mejor
  caso (arreglo ya ordenado) y $\mathcal{O}(n \log n)$ en el peor.
  \par\smallskip
  \textit{Referencia:} Peters, T. (2002). ``[Python-Dev] Sorting''.
  \url{https://bugs.python.org/issue4255}\\
  Artículo descriptivo: \url{https://en.wikipedia.org/wiki/Timsort}
\end{tecnico}

\begin{ejemplo}[MergeSort en la vida real: \texttt{git merge}]
  El comando \texttt{git merge} que une ramas de código tiene una lógica análoga
  a MergeSort: toma dos secuencias de cambios (commits) ya ordenadas por fecha
  y las fusiona en una sola historia coherente, resolviendo conflictos cuando
  los mismos archivos fueron modificados en ambas ramas. No es el mismo algoritmo,
  pero el patrón de ``fusionar dos secuencias ordenadas'' está en el corazón de
  cómo Git maneja la historia del código.
\end{ejemplo}

% =============================================================================
\section{Ejercicios}
% =============================================================================

\subsection*{Ejercicio 1 — Traza manual}

Aplica MergeSort al arreglo \texttt{[5, 3, 8, 1, 4, 2, 7, 6]} a mano (lápiz
y papel). Dibuja el árbol de divisiones y la secuencia de fusiones, mostrando
el estado del arreglo después de cada llamada a \cpp{merge}.

\subsection*{Ejercicio 2 — Contar comparaciones}

Modifica la función \cpp{merge} para que cuente el número de comparaciones
(\cpp{arr[i] <= arr[j]}) realizadas. Luego modifica \cpp{mergesort} para
acumular ese conteo y al final imprime cuántas comparaciones se hicieron en
total para el arreglo de 8 elementos.

\noindent¿Coincide con la predicción teórica de $n \log_2 n = 8 \times 3 = 24$?

\subsection*{Ejercicio 3 — Orden descendente}

Modifica \cpp{merge} para que el resultado quede en orden \textbf{descendente}.
¿Cuántas líneas de código cambiaste? ¿Qué te dice eso sobre el diseño del algoritmo?

\subsection*{Ejercicio 4 — MergeSort genérico con comparador}

Modifica las firmas de \cpp{merge} y \cpp{mergesort} para que acepten un
parámetro \cpp{bool(*cmp)(int, int)} — un apuntador a función comparadora,
igual que hicimos con Bubble Sort. Verifica que con \cpp{cmp = [](int a, int b)
\{ return a <= b; \}} el resultado es ascendente y con \cpp{a >= b} es
descendente, sin cambiar el cuerpo de las funciones.

% =============================================================================
\section{Resumen}
% =============================================================================

\begin{recuerda}
\begin{itemize}
  \item MergeSort divide el arreglo hasta llegar a subarreglos de 1 elemento
        (\textbf{caso base}), luego los fusiona de abajo hacia arriba.
  \item \textbf{Complejidad}: $\mathcal{O}(n \log n)$ en todos los casos —
        el mejor, el promedio y el peor.
  \item \textbf{Memoria extra}: $\mathcal{O}(n)$ para el buffer temporal en
        \cpp{merge}.
  \item \textbf{Estable}: preserva el orden relativo de elementos iguales.
  \item Es \textbf{óptimo}: ningún algoritmo de comparaciones puede hacerlo
        mejor (prueba del límite inferior $\Omega(n \log n)$).
  \item QuickSort es más rápido en la práctica (cache), pero $\mathcal{O}(n^2)$
        en el peor caso. MergeSort garantiza $\mathcal{O}(n \log n)$ siempre.
  \item TimSort (Python, Java, \cpp{std::stable\_sort}) es MergeSort adaptado
        para arreglos reales con runs ya ordenados.
  \item \cpp{int* arr} en la firma significa \textit{paso por referencia}:
        la función modifica el arreglo original sin copiarlo.
\end{itemize}
\end{recuerda}

\end{document}
